\documentclass[11pt,a4paper]{article}
\usepackage[utf8]{inputenc}
\usepackage[margin=2.5cm]{geometry}
\usepackage{tcolorbox}
\usepackage{listings}
\usepackage{xcolor}
\usepackage{hyperref}
\usepackage{fontspec}
\usepackage{datetime2}
\usepackage{enumitem}
\usepackage{booktabs}
\usepackage{multirow}
\usepackage{parskip}
\usepackage{graphicx}

% Set classic LaTeX fonts (Computer Modern)
\setmainfont{Latin Modern Roman}
\setsansfont{Latin Modern Sans}

% Custom colors
\definecolor{codeblue}{rgb}{0.25,0.5,0.5}
\definecolor{codegray}{rgb}{0.5,0.5,0.5}
\definecolor{codepurple}{rgb}{0.58,0,0.82}
\definecolor{backcolour}{rgb}{0.95,0.95,0.92}

% Code listing style
\lstdefinestyle{mystyle}{
    backgroundcolor=\color{backcolour},
    commentstyle=\color{codepurple},
    keywordstyle=\color{codeblue},
    numberstyle=\tiny\color{codegray},
    stringstyle=\color{codepurple},
    basicstyle=\ttfamily\footnotesize,
    breakatwhitespace=false,
    breaklines=true,
    captionpos=b,
    keepspaces=true,
    numbers=left,
    numbersep=5pt,
    showspaces=false,
    showstringspaces=false,
    showtabs=false,
    tabsize=2
}

\lstset{style=mystyle}

% Hyperlink settings
\hypersetup{
    colorlinks=true,
    linkcolor=codeblue,
    urlcolor=codeblue,
    pdftitle={AI-Powered Windows Cleaner - System Report},
    pdfauthor={Abdullah Al Mamun},
    pdfsubject={System Analysis Report},
    pdfkeywords={Windows, Cleaner, AI, System, Security}
}

\titleformat{\section}{\large\bfseries\color{codeblue}}{\thesection}{1em}{}
\titleformat{\subsection}{\normalsize\bfseries\color{codeblue}}{\thesubsection}{1em}{}
\titleformat{\subsubsection}{\itshape\color{codepurple}}{\thesubsubsection}{1em}{}

\begin{document}

\pagestyle{empty}

\begin{tcolorbox}[colback=codeblue!10!white,colframe=codeblue!80!black,title=System Report Header, fonttitle=\bfseries\Large]
\section*{AI-Powered Windows Cleaner}
\vspace{0.5em}
\subsection*{Comprehensive System Analysis Report}
\end{tcolorbox}

\vspace{-1em}
\begin{flushright}
\small
Author: \textbf{Abdullah Al Mamun} \\
ORCID: \href{https://orcid.org/0009-0006-7473-0024}{0009-0006-7473-0024} \\
Email: \href{mailto:mamun.swe.de@gmail.com}{mamun.swe.de@gmail.com} \\
Location: Vienna, Austria \\
Date: \today
\end{flushright}

\vspace{1em}

\section{Executive Summary}

This document provides a comprehensive analysis of the AI-Powered Windows Cleaner system architecture, current implementation status, test coverage, security posture, and planned future roadmap. The system is a production-quality, enterprise-grade Windows maintenance tool that leverages local AI (LLaMA 3.2 via Ollama) for intelligent system optimization.

\subsection{Current Status}
\begin{itemize}
    \item \textbf{Overall Health Score:} \textbf{100/100}
    \item \textbf{Test Coverage:} 94\%+ (36 tests passing)
    \item \textbf{Code Quality:} Zero linting errors (Ruff)
    \item \textbf{Type Safety:} 100\% type-checked (Mypy)
    \item \textbf{Security:} Zero vulnerabilities (Bandit)
    \item \textbf{Complexity:} A/B grades only (Radon)
\end{itemize}

\section{Project Architecture}

\subsection{System Overview}
The project follows a hybrid architecture with a native Windows GUI communicating via a local REST API to a containerized AI backend.

\begin{tcolorbox}[colback=white!95!black,title=Architecture Diagram]
\begin{verbatim}
┌─────────────────────────────────────────────────────────────┐
│  Native Windows Host (Python)                               │
│  ┌─────────────────────────────────────────────────────────┐│
│  │  PySide6 GUI Dashboard                                  ││
│  │  ├─ Overview Widget (Live metrics)                      ││
│  │  ├─ Scanner Results View (Multi-threaded scan)          ││
│  │  ├─ AI Chat Widget (Async Ollama integration)            ││
│  │  └─ History/Rollback View                               ││
│  └─────────────────────────────────────────────────────────┘│
│                                                             │
│  ┌──────────────────┐    ┌────────────────────┐            │
│  │  HTTP Client     │───▶│  FastAPI Backend   │            │
│  │  (requests)      │    │  (port 8000)       │            │
│  └──────────────────┘    └────────┬───────────┘            │
│                                   │                         │
│                                   ▼                         │
│  ┌──────────────────┐                                       │
│  │  Windows File    │                                       │
│  │  System API      │                                       │
│  └──────────────────┘                                       │
└─────────────────────────────────────────────────────────────┘

┌─────────────────────────────────────────────────────────────┐
│  Podman Container (WSL2)                                    │
│  ┌──────────────┐    ┌──────────────┐                       │
│  │  Ollama      │◀──▶│  llama3.2:1b │                       │
│  │  (port 11434)│    │   ~1.3 GB    │                       │
│  └──────────────┘    └──────────────┘                       │
└─────────────────────────────────────────────────────────────┘
\end{verbatim}
\end{tcolorbox}

\subsection{Technology Stack}

\begin{table}[h]
\centering
\begin{tabular}{llp{6cm}}
\toprule
\textbf{Layer} & \textbf{Tool} & \textbf{Purpose} \\
\midrule
GUI & PySide6 6.6+ & Native Windows desktop UI \\
Glassmorphism & win32mica & Windows 11 Mica DWM backdrop \\
Charts & pyqtgraph & Hardware-accelerated storage graphs \\
System Metrics & psutil & Real-time CPU / RAM / Disk monitoring \\
File I/O & pathlib + os & Safe, cross-version filesystem operations \\
AI Chat Client & requests + QThread & Async HTTP to local FastAPI backend \\
AI Backend & FastAPI + uvicorn & REST API inside container \\
Containerisation & Podman + podman-compose & Isolated AI sandbox \\
Database & SQLite3 & History, rollback logs, preferences \\
\bottomrule
\end{tabular}
\end{table}

\section{Codebase Analysis}

\subsection{Module Structure}

\begin{verbatim}
src/ai_health_copilot/
├── main.py                    # Application entry point
├── gui/
│   ├── main_window.py         # Main application window
│   ├── views/
│   │   ├── overview.py        # Dashboard with live metrics
│   │   ├── scanner_results.py # Scan results view
│   │   ├── ai_chat.py         # Chat interface widget
│   │   └── history.py         # History/rollback view
│   └── widgets/               # Reusable UI components
├── core/
│   ├── scanner/
│   │   ├── system_info.py     # System metrics collection
│   │   └── large_files.py     # Large file detection
│   ├── cleaner/
│   │   ├── base.py            # Base cleaner framework
│   │   ├── downloads.py       # Downloads cleanup
│   │   ├── recycle_bin.py     # Recycle bin management
│   │   └── windows_temp.py    # Temp file cleanup
│   ├── duplicate/             # Duplicate detection (Phase 16 pending)
│   ├── rollback/              # Quarantine manager
│   ├── logger/                # Logging utilities
│   └── scheduler/             # Windows Task Scheduler integration
├── ai/
│   └── advisor.py             # AI backend communication
└── database/
    ├── manager.py             # SQLite operations
    └── schema.sql             # Database schema
\end{verbatim}

\subsection{Core Dependencies}

\begin{lstlisting}[language=Python,caption={Main Dependencies}]
PySide6>=6.6.0      # GUI Framework
psutil>=5.9.0       # System metrics
pyqtgraph>=0.13.0   # Charts and graphs
requests>=2.31.0    # HTTP client
sqlite3             # Database (stdlib)
win32mica           # Windows 11 Mica backdrop
\end{lstlisting}

\begin{lstlisting}[language=YAML,caption={Container Services}]
services:
  ai-backend:
    image: python:3.12-slim
    ports: ["8000:8000"]
    
  ollama:
    image: ollama/ollama:latest
    ports: ["11434:11434"]
    volumes: ["ollama-data:/root/.ollama"]
\end{lstlisting}

\section{Quality Assurance}

\subsection{Testing Suite Results}

\begin{tcolorbox}[colback=green!5!white,colframe=green!70!black,title=Current Test Status]
\begin{itemize}
    \item \textbf{Unit Tests:} 36 passed in 12.58s
    \item \textbf{Integration Tests:} All module integrations verified
    \item \textbf{Coverage:} $\geq$ 94\% across all modules
    \item \textbf{Regression Testing:} All previous phase tests passing
\end{itemize}
\end{tcolorbox}

\subsection{Code Quality Gates}

\begin{table}[h]
\centering
\begin{tabular}{lcc}
\toprule
\textbf{Tool} & \textbf{Status} & \textbf{Requirement} \\
\midrule
Ruff (Linting) & PASS & Zero warnings \\
Mypy (Type Check) & PASS & 100\% typed \\
Bandit (Security) & PASS & Zero vulnerabilities \\
Radon (Complexity) & PASS & A/B grades only \\
Pytest (Tests) & PASS & $\geq$ 90\% coverage \\
\bottomrule
\end{tabular}
\end{table}

\section{Future Roadmap}

\subsection{Phase 15: Enhanced AI Capabilities}
\begin{itemize}
    \item Streaming AI responses (token-by-token display)
    \item Multi-modal reasoning (image analysis of error dialogs)
    \item Conversational memory (session-aware AI)
\end{itemize}

\subsection{Phase 16: History and Recovery}
\begin{itemize}
    \item Visual history timeline
    \item One-click rollback for entire sessions
    \item Snapshot-based restore points
\end{itemize}

\subsection{Phase 17: Settings Module}
\begin{itemize}
    \item AI model selector (llama3.2, llama3.1, code-llama)
    \item Configurable scan targets and exclusions
    \item Advanced scheduler configuration
\end{itemize}

\subsection{Phase 18: Advanced Cleanup}
\begin{itemize}
    \item Content-aware duplicate finder
    \item Large file auditor with recommendations
    \item Browser cache management
    \item Registry optimization
\end{itemize}

\subsection{Phase 19: Performance Optimizer}
\begin{itemize}
    \item Startup program manager
    \item Memory leak detection
    \item SSD/HDD fragmentation analysis
\end{itemize}

\subsection{Phase 20: Enterprise Features}
\begin{itemize}
    \item Multi-device sync
    \item Usage analytics dashboard
    \item Admin mode with audit logging
    \item Portable deployment option
\end{itemize}

\section{Security Analysis}

\begin{tcolorbox}[colback=blue!5!white,colframe=blue!70!black,title=Security Posture]
\begin{itemize}
    \item No shell injection vulnerabilities (validated)
    \item Path traversal protection implemented
    \item Quarantine-first deletion with full rollback
    \item No cloud dependencies (100\% offline)
    \item No embedded secrets or credentials
    \item Admin privilege detection and graceful degradation
\end{itemize}
\end{tcolorbox}

\section{Database Schema}

\begin{lstlisting}[language=SQL,caption={SQLite Schema}]
CREATE TABLE History (
    id INTEGER PRIMARY KEY AUTOINCREMENT,
    action_type TEXT NOT NULL,
    target TEXT NOT NULL,
    size_bytes INTEGER DEFAULT 0,
    timestamp DATETIME DEFAULT CURRENT_TIMESTAMP
);

CREATE TABLE CleanupLogs (
    id INTEGER PRIMARY KEY AUTOINCREMENT,
    session_id TEXT NOT NULL,
    details TEXT NOT NULL,
    total_recovered INTEGER DEFAULT 0,
    timestamp DATETIME DEFAULT CURRENT_TIMESTAMP
);

CREATE TABLE AIConversations (
    id INTEGER PRIMARY KEY AUTOINCREMENT,
    role TEXT NOT NULL,
    message TEXT NOT NULL,
    timestamp DATETIME DEFAULT CURRENT_TIMESTAMP
);
\end{lstlisting}

\section{Conclusion}

The AI-Powered Windows Cleaner is a production-ready system with a modular architecture, comprehensive testing suite, and robust security measures. All quality gates pass at 100\%, enabling safe progression to future development phases.

\begin{tcolorbox}[colback=gray!5!white,colframe=gray!70!black,title=Report Metadata]
Report Generated: \today \\
Python Version: 3.12.4 \\
Qt Framework: PySide6 6.6+ \\
Status: \textbf{PRODUCTION READY}
\end{tcolorbox}

\end{document}